% Source: https://artofproblemsolving.com/community/c3h2002937_exponential_function_proof
% Stable identity: aops-topic-2002937-q01
\textbf{(a)} For every real $r$ we have $r^2\ge0>b$. So no real number squares to $b$. Since the rules of exponents force $\left(b^{\dfrac{1}{2}}\right)^2=b^1=b$, the expression $b^{\dfrac{1}{2}}$ has no real value. Hence $f(x)=a\cdot b^x$ is undefined at $x=\dfrac{1}{2}$ (and similarly at every $x=\dfrac {p}{q}$ in lowest terms with $q$ even).

\textbf{(b)} Real cube roots of negative numbers exist: $(-8)^{\dfrac{1}{3}}=\sqrt[3]{-8}=\boxed{-2}$ and $(-8)^{\dfrac{2}{3}}=\left(\sqrt[3]{-8}\right)^2=\boxed{4}$. But by (a), $(-8)^{\dfrac{1}{2}}$ is not real. So $x\mapsto(-8)^x$ is defined only on some rational numbers (those with odd denominator), and even there it jumps between negative and positive values. It cannot be a real-valued function on all of $\mathbb R$, which is why the definition requires $b>0$. $\blacksquare$
